% page 43 du fac-simile (image p-042 du scan)
% bloc 175.3 x 237.7 mm ; marge gauche 26.7 mm
\pgc{5.080mm}{0.000mm}{
\l{\cel{58}35}
\l{}
\l{}
\l{\sou{apostilo}. Noto sur peticiono, por rekomendar lu. - EFIS.}
\l{}
\l{\sou{apostolo}. (\sou{religio katol}.) Singla ek la dek-e-du dicipuli qui}
\l{\cel{3}recevis de Iesu-Kristo la komiso predikor evangelio. -}
\l{\cel{3}Ta qua propagas la doktrino Kristana. - DEFIR.}
\l{}
\l{\sou{apostrofo}.\cel{2}I. (\sou{gram}.) Signo ortografiala qua remplasas litero}
\l{\cel{3}elizionita. - II. (\sou{retoriko}). Procedo oratorala per qua onu}
\l{\cel{3}interpelas subite persono prezenta, ento nevidebla, kozo quan}
\l{\cel{3}onu traktas quale se lu esus persono. - Interpelo vigoroza}
\l{\cel{3}di ulu. - DEFIRS.}
\l{}
\l{\sou{apoteko}.\cel{2}Butiko ube la medikamenti preparesas, konservesas,vendesas.}
\l{\cel{3}a konsumeriv-DEFIRS.}
\l{}
\l{\sou{apotemo}. (\sou{geom}.) I. Radio di la cirklo enskribata en poligono re-}
\l{\cel{3}gulala. - II. Alteso di un ek la edri di piramido regulala. -}
\l{\cel{3}DEFIRS.}
\l{}
\l{\sou{apoteoso}. I. Granto di honori deala. - II. Granto di honori}
\l{\cel{3}extraordinara. - DEFIRS.}
\l{}
\l{\sou{apoziciono}.\cel{2}(\sou{gram}.) Qualifiko quan onu adjuntas a substantivo}
\l{\cel{3}per altra substantivo apud-pozata. - DEFIRS.}
\l{}
\l{\sou{apro}.\cel{2}(\sou{zool}.) Mamifero, ek la pakidermi, ek genero qua vivas}
\l{\cel{3}sovaja, ed egardesas kom la tronko di la porko domestika. -}
\l{\cel{3}L.\cel{1}\sou{aper}.\cel{2}DRL.}
\l{}
\l{\sou{aprentiso}. (\sou{pri}) Ta qua lernas mestiero. - DEFIS.}
\l{}
\l{\sou{apretar}. (\sou{trans.}) (\sou{tekn.}) I. Submisar ula texuri a preparajo qua}
\l{\cel{3}igas li koheranta, solida, glata, briloza. - II. Indutar telo}
\l{\cel{3}sur qua onu esas piktonta. - DFIR.}
\l{}
\l{\sou{aprilo.}\cel{2}La monato quaresma di la kalendar-yaro. - DEFIRS.}
\l{}
\l{\sou{a priori\textquotedbl{}}.\cel{2}Segun donaji qui pre-iras experimentado. -DEFIRS.}
\l{}
\l{\sou{aprobar}.\cel{2}(\sou{trans.}) I. Judikar bona, laudinda. - II. Deklarar}
\l{\cel{3}(per formulo) ke akto, konto, esas exakta. - DEFIRS.}
\l{}
\l{\sou{aprocho}. (\sou{fortifikuro}). Terlanoruro per qua militisti +asiejanta}
\l{\cel{3}avancas a la fortifikuri e rempari enemikala. - DEFIRS.}
\l{}
\l{\sou{aproximar}. (\sou{trans.}) (\sou{tekn.}) Evaluar tale ke onu saciesas per}
\l{\cel{3}nura su-proximigo a lo exakta.}
\l{}
\l{\sou{apsido}. I. (\sou{arkeol}.) Mi-cirklo qua formacas la fundo di la bazi-}
\l{\cel{3}liki pagana, e divenis la santuario di la baziliki Kristana.}
\l{\cel{3}- II. (\sou{astron}.) Punto di la orbito di planeto ube lu distas}
\l{\cel{3}maxime o minime de la astro cirkum qua lu su movas. - DEFIRS.}
}
